\documentclass[12pt]{article}
\input{article_style.tex}

\title{Raport projektowy\\Uczenie się ze wzmocnieniem}
\author{Igor Paśko}
\date{2026/2027}

\begin{document}
\maketitle
\tableofcontents
\newpage

\section{Informacje podstawowe}

\begin{center}
\begin{tabular}{p{0.34\textwidth}p{0.56\textwidth}}
\toprule
Element & Informacja \\
\midrule
Imię i nazwisko & Igor Paśko \\
Numer albumu & 177137 \\
Kierunek/specjalność & 3EF-DI/AI \\
Tytuł projektu & Sterowanie agentem w grze unikania pocisków \\
Link do repozytorium & \url{https://github.com/Igor-przinf/RL-bullet-dodge} \\
Data oddania & Do uzupełnienia \\
\bottomrule
\end{tabular}
\end{center}

\section{Temat i cel projektu}

\subsection{Opis problemu}
Rozpatrywanym problemem jest sterowanie obiektem (agentem) na dwuwymiarowej planszy w czasie rzeczywistym, którego głównym celem jest unikanie nadciągających z różnych kierunków pocisków (przeszkód) i przetrwanie jak najdłuższego czasu.

\subsection{Cel agenta}
Celem agenta jest opanowanie polityki ruchu, która pozwala mu unikać kolizji z pociskami przez maksymalny czas trwania epizodu. Działanie agenta uznaje się za skuteczne, gdy w procesie uczenia średni czas przeżycia oraz liczba unikniętych pocisków systematycznie rosną w stosunku do agenta bazowego.

\subsection{Uzasadnienie zastosowania uczenia się ze wzmocnieniem}
Zastosowanie uczenia się ze wzmocnieniem (RL) w tym problemie jest w pełni uzasadnione ze względu na sekwencyjny charakter podejmowanych decyzji:
\begin{itemize}
    \item Agent podejmuje decyzje (ruch w określonym kierunku) w każdym kroku czasowym, co bezpośrednio wpływa na jego pozycję na planszy.
    \item Decyzje podjęte w danej chwili determinują przyszły stan środowiska oraz dostępną przestrzeń do manewru unikowego przed kolejnymi pociskami.
    \item Istotny jest wynik długoterminowy (całkowity czas przeżycia w epizodzie), a nie tylko natychmiastowa nagroda za pojedynczy unik.
    \item Problem nie ma charakteru statycznej klasyfikacji, lecz wymaga dynamicznej adaptacji strategii w interakcji ze zmieniającym się otoczeniem.
\end{itemize}

\section{Formalizacja problemu RL/MDP}

\subsection{Agent i środowisko}
Agentem jest obiekt poruszający się po dwuwymiarowej planszy. Środowisko generuje pociski pojawiające się na krawędziach ekranu, które przemieszczają się z określoną prędkością. W pojedynczym kroku epizodu agent wykonuje akcję ruchu, środowisko aktualizuje pozycje pocisków, sprawdza kolizje i zwraca nowy stan oraz nagrodę.

\subsection{Stan lub obserwacja}
Wektor stanu w kroku czasowym $t$ opisuje pozycję agenta oraz najbliższe pociski:

\begin{center}
\begin{tabular}{p{0.30\textwidth}p{0.60\textwidth}}
\toprule
Składowa stanu/obserwacji & Znaczenie \\
\midrule
$x_a, y_a$ & Współrzędne pozycji agenta na planszy \\
$\Delta x_i, \Delta y_i$ & Wektor względnej odległości do $i$-tego pocisku \\
$v_{xi}, v_{yi}$ & Składowe prędkości $i$-tego pocisku \\
\bottomrule
\end{tabular}
\end{center}

\subsection{Akcje}
Agent dysponuje dyskretną przestrzenią akcji pozwalającą na sterowanie na płaszczyźnie:

\begin{center}
\begin{tabular}{p{0.30\textwidth}p{0.60\textwidth}}
\toprule
Akcja & Opis \\
\midrule
0 & Brak ruchu (pozostanie w miejscu) \\
1 & Ruch w górę \\
2 & Ruch w dół \\
3 & Ruch w lewo \\
4 & Ruch w prawo \\
\bottomrule
\end{tabular}
\end{center}

\subsection{Funkcja nagrody}
Nagroda motywuje agenta do przetrwania i unikania zagrożeń:

\begin{equation}
    r_t = 
\begin{cases} 
-100 & \text{kolizja z pociskiem (koniec epizodu)}, \\
+1.0 & \text{za każdy przetrwany krok czasowy}, \\
-0.1 & \text{kara za zbliżenie się do krawędzi planszy}.
\end{cases}
\end{equation}

W tej sekcji wyjaśniono przyjęte założenia:
\begin{itemize}
    \item Agent otrzymuje nagrodę dodatnią (+1.0) za każdy krok, w którym nie doszło do kolizji, co premiuje długie przeżycie.
    \item Kara (-100) jest nakładana natychmiastowo przy kolizji z pociskiem, co kończy epizod.
    \item Występuje niewielka kara za krawędzie, aby zapobiec utknięciu agenta w rogu planszy.
\end{itemize}

\subsection{Warunki zakończenia epizodu}
Epizod kończy się w przypadku:
\begin{enumerate}
    \item Kolizji agenta z dowolnym pociskiem (porażka).
    \item Osiągnięcia maksymalnego limitu kroków w epizodzie, np. $T_{max} = 1000$ (sukces).
\end{enumerate}

\subsection{Założenia i ograniczenia środowiska}
Środowisko zakłada uproszczoną fizykę 2D (ruch prostoliniowy pocisków, dyskretne kroki czasowe). Pomięto zaawansowane efekty wizualne czy rotację pocisków, skupiając się wyłącznie na relacjach przestrzennych i decyzjach unikowych.

\section{Implementacja}

\subsection{Struktura repozytorium}
Projekt będzie opierał się na czystej strukturze katalogów przechowującej kod środowiska, skrypty trenujące oraz raport w LaTeX-u:

\begin{verbatim}
repozytorium/
|-- src/
|   |-- environment.py
|   |-- agent.py
|   `-- main.py
|-- experiments/
|-- report/
|-- results/
`-- README.md
\end{verbatim}

\subsection{Środowisko}
Środowisko zostanie zaimplementowane w języku Python z wykorzystaniem biblioteki \texttt{Pygame} lub jako dedykowana klasa kompatybilna z interfejsem \texttt{Gymnasium}. Posiadać będzie funkcję \texttt{reset()} inicjalizującą pozycję gracza i listę pocisków oraz funkcję \texttt{step(action)}, która wykonuje ruch, aktualizuje pozycje pocisków, wykrywa kolizje i zwraca krotkę \texttt{(next\_state, reward, done, info)}.

\subsection{Agent bazowy}
Jako punkt odniesienia (baseline) zostanie zaimplementowany agent losowy (\textit{Random Agent}), który w każdym kroku wybiera akcję losowo z rozkładu jednostajnego. Posłuży on do wykazania, że algorytm uczenia ze wzmocnieniem faktycznie uczy się efektywnej polityki unikania.

\subsection{Agent uczący się}

Należy opisać zastosowany algorytm RL lub DRL, jego parametry oraz sposób implementacji. W przypadku Q-tab należy opisać reprezentację tablicy wartości. W przypadku DRL należy opisać wejścia modelu, wyjścia modelu, architekturę oraz sposób uczenia.

\section{Plan eksperymentów}

Należy opisać eksperymenty zaplanowane w celu oceny działania agenta.

\begin{center}
\begin{tabular}{p{0.24\textwidth}p{0.30\textwidth}p{0.30\textwidth}}
\toprule
Eksperyment & Co jest zmieniane? & Co jest mierzone? \\
\midrule
E1 & Do uzupełnienia & Do uzupełnienia \\
E2 & Do uzupełnienia & Do uzupełnienia \\
E3 & Do uzupełnienia & Do uzupełnienia \\
\bottomrule
\end{tabular}
\end{center}

W tej sekcji należy wskazać:
\begin{itemize}
    \item liczbę epizodów uczenia,
    \item liczbę powtórzeń eksperymentu albo uzasadnienie braku powtórzeń,
    \item parametry uczenia,
    \item metryki oceny,
    \item warianty środowiska, map lub funkcji nagrody,
    \item sposób porównania z agentem bazowym.
\end{itemize}

\section{Wyniki}

W tej sekcji należy przedstawić wyniki eksperymentów. Wyniki powinny być pokazane w postaci tabel, wykresów lub czytelnych zestawień.

\begin{center}
\begin{tabular}{p{0.28\textwidth}p{0.20\textwidth}p{0.20\textwidth}p{0.20\textwidth}}
\toprule
Wariant & Średni zwrot & Skuteczność & Długość epizodu \\
\midrule
Agent bazowy & Do uzupełnienia & Do uzupełnienia & Do uzupełnienia \\
Agent RL/DRL & Do uzupełnienia & Do uzupełnienia & Do uzupełnienia \\
\bottomrule
\end{tabular}
\end{center}

\begin{figure}[h]
    \centering
    \fbox{\parbox{0.82\textwidth}{\centering Miejsce na wykres procesu uczenia}}
    \caption{Przykładowy wykres procesu uczenia agenta.}
\end{figure}

\section{Analiza i wnioski}

Należy zinterpretować wyniki. W tej części należy odpowiedzieć m.in. na pytania:
\begin{itemize}
    \item Czy agent uczący się działa lepiej niż agent bazowy?
    \item Jak zmieniało się zachowanie agenta w trakcie uczenia?
    \item Które elementy środowiska były najtrudniejsze?
    \item Czy funkcja nagrody prowadziła do oczekiwanego zachowania?
    \item Czy wystąpiły przypadki niepożądanego zachowania agenta?
    \item Jakie są najważniejsze ograniczenia rozwiązania?
\end{itemize}

\section{Instrukcja uruchomienia}

Należy opisać sposób uruchomienia projektu od czystego repozytorium. Przykładowo:

\begin{verbatim}
git clone https://github.com/...
cd nazwa-repozytorium
python -m venv .venv
.venv\Scripts\activate
pip install -r requirements.txt
python src/main.py
\end{verbatim}

Należy również wskazać, jak odtworzyć eksperymenty oraz gdzie zapisywane są wyniki.

\section{Wsparcie AI}

W tej sekcji należy obowiązkowo opisać wykorzystanie narzędzi sztucznej inteligencji. Jeżeli student nie korzystał z takich narzędzi, należy wpisać jednoznaczną informację, że nie korzystano ze wsparcia AI.

\begin{center}
\begin{tabular}{p{0.13\textwidth}p{0.19\textwidth}p{0.21\textwidth}p{0.27\textwidth}}
\toprule
Data & Narzędzie/model & Cel użycia & Transkrypt / link / zakres \\
\midrule
Do uzupełnienia & Do uzupełnienia & Do uzupełnienia & Do uzupełnienia \\
Do uzupełnienia & Do uzupełnienia & Do uzupełnienia & Do uzupełnienia \\
\bottomrule
\end{tabular}
\end{center}

Należy zapisać transkrypty rozmów z modelami AI wraz z datą oraz nazwą modelu. Jeżeli rozmowa odbywała się poza narzędziami programistycznymi, należy dodatkowo udostępnić całą rozmowę poprzez link, o ile użyte narzędzie umożliwia takie udostępnienie.

Student powinien także wyjaśnić, które elementy wygenerowane lub zasugerowane przez AI zostały samodzielnie sprawdzone, zmodyfikowane albo odrzucone.

\section{Literatura}

W tej sekcji należy wskazać wykorzystane źródła: książki, artykuły naukowe, dokumentację bibliotek, materiały wykładowe oraz inne źródła użyte podczas pracy nad projektem.

\begin{thebibliography}{9}
\bibitem{sutton_barto}
R. S. Sutton, A. G. Barto, \textit{Reinforcement Learning: An Introduction}, 2nd ed., MIT Press, 2018.
\end{thebibliography}

\end{document}
